(1) 定义 $\omega_c=qB/m$。最小耦合 $H=(\mathbf p-q\mathbf A)^2/(2m)+m\omega_0^2r^2/2$ 展开为 $H=p^2/(2m)+\frac m2(\omega_0^2+\omega_c^2/4)r^2-(\omega_c/2)L_z$，因为 $\nabla\cdot\mathbf A=0$。令 $\Omega=\sqrt{\omega_0^2+\omega_c^2/4}$，得到二维频率 $\Omega$ 的各向同性谐振子加角动量劈裂。
(2) 二维谐振子径向本征能为 $\hbar\Omega(2n+|l|+1)$，而 $L_z$ 本征值为 $\hbar l$。因此完整能谱 $E_{n,l}=\hbar\Omega(2n+|l|+1)-\hbar\omega_c l/2$；当 $l>0$ 时轨道磁矩项降低能量，但 $\Omega>\omega_c/2$ 确保有限 $\omega_0$ 下能级下有界。
(3) 基态 $n=l=0$ 的归一化波函数 $\psi_{0,0}(r,\phi)=(\pi a_\Omega^2)^{-1/2}\exp[-r^2/(2a_\Omega^2)]$，$a_\Omega=\sqrt{\hbar/(m\Omega)}$。能量 $E_{0,0}=\hbar\Omega$，且由极坐标积分 $\langle r^2\rangle=2\pi\int_0^\infty r^3|\psi|^2dr=a_\Omega^2$，磁场增大使基态径向宽度下降。
(4) $B\to0$ 时 $\omega_c\to0$、$\Omega\to\omega_0$，恢复二维谐振子的角动量简并能谱。强磁场时 $\Omega\simeq\omega_c/2+\omega_0^2/\omega_c$，对固定 $n=0,l\ge0$，能级趋 $\hbar\omega_c/2$，对应最低Landau级的近简并；有限 $\omega_0$ 仍引入随 $l$ 增加的小正劈裂，而严格无束缚极限 $\omega_0=0$ 恢复Landau简并结构。
